% !PRINTED 9
\weberpage{9}

Die Nebengruppen haben nicht die Merkmale einer Gruppe. Denn damit zwei
Elemente aus $Q_1$ bei der Zusammensetzung wieder ein Element von $Q_1$
ergeben, müsste etwa
\[
a_1 b a_2 b = a_3 b,
\]
also $a_1 b a_2 = a_3$, $b = a_1^{-1} a_3 a_2^{-1}$ sein. Es wäre also $b$ der
Voraussetzung entgegen in $Q$ enthalten.

Die Benennung Nebengruppe ist also nur uneigentlich zu verstehen.

Zwei Theiler $Q$, $Q'$ von $P$ haben immer das Element 1 mit einander
gemein. Sie können aber auch noch andere Elemente gemeinschaftlich haben,
und diese gemeinschaftlichen Elemente bilden eine Gruppe, die wir mit $R$
bezeichnen wollen. Denn gehören die Elemente $a$ und $b$ sowohl zu $Q$ als
zu $Q'$, so gilt wegen des Gruppencharakters von $Q$ und $Q'$ dasselbe von
dem Compositum $ab$; also ist $R$ eine Gruppe. Diese gemeinschaftliche
Gruppe nennen wir den \gesperrt{grössten gemeinschaftlichen Theiler} von $Q$ und $Q'$
oder auch, mit einem geometrischen Anklange, den \gesperrt{Durchschnitt von} $Q$ und
$Q'$. Ebenso folgt, dass die Elemente, die irgend einer Anzahl von
Theilern von $P$ gemeinsam sind, eine Gruppe bilden, die wir ebenso als
den grössten gemeinschaftlichen Theiler oder den Durchschnitt aller dieser
Gruppen bezeichnen.

In jeder Gruppe können wir nach folgendem Verfahren Theiler bilden.

Ein Theiler ist immer das Einheitselement für sich.

Bezeichnen wir die wiederholte Zusammensetzung eines Elementes mit sich
selbst durch Potenzen (mit $a^0$ das Einheitselement), so kann die Reihe
der Elemente
\[
1,\ a,\ a^2,\ a^3 \ldots,
\]
die alle der Gruppe $P$ angehören, nicht lauter verschiedene Elemente
enthalten. Ist also das erste Element, das zum zweiten Male wiederkehrt,
\[
a^\mu = a^{\mu+\alpha},
\]
so folgt, dass $a^\alpha = 1$ sein muss, und dass $\alpha$ die kleinste
positive Zahl ist, die dieser Bedingung genügt. Es ist dann, wenn
$m = q\alpha$ ein beliebiges Vielfaches von $\alpha$ ist, $a^m = 1$ und
umgekehrt: so oft $a^m = 1$ ist, muss $m$ durch $\alpha$ theilbar sein;
denn sonst könnte man $m = q\alpha + \alpha'$ setzen, worin $\alpha'$
positiv und kleiner als $\alpha$ ist, und es wäre $a^{\alpha'} = 1$, gegen
die Voraussetzung,
%CONT>>
